\documentclass[11pt,a4paper]{article}
\usepackage[T1]{fontenc}
\usepackage{lmodern,amsmath,amssymb,amsthm,microtype}
\usepackage[margin=27mm]{geometry}
\usepackage[colorlinks=true,linkcolor=blue,citecolor=blue,urlcolor=blue]{hyperref}
\setlength{\emergencystretch}{2em}
\newtheorem{theorem}{Theorem}[section]
\newtheorem{lemma}[theorem]{Lemma}
\newtheorem{corollary}[theorem]{Corollary}
\newtheorem{proposition}[theorem]{Proposition}
\newcommand{\Q}{\mathbb Q}
\DeclareMathOperator{\Tr}{Tr}
\title{Vertical valuations in abelian families of actual traces}
\author{}\date{17 November 2026}
\begin{document}\maketitle
\begin{abstract}
A fixed good prime of positive odd weight valuation obstructs a
nonconstant abelian family at a fixed slope in the actual-trace
detectors. A quotient by the connected kernel of the point map
preserves a putative square root and has good reduction at that prime.
The generic vertical valuation of a nonconstant reciprocal abscissa
then forces the radicand to have odd valuation. Combined with the
nonconstant-slope obstruction, this excludes nonconstant rational
abelian families through prepared points in the fixed-root derivative
criterion and the fixed-input general criterion. Individual rational
witnesses remain possible. A separate argument allows the derivative
root point to vary while the slope depends rationally and affinely on
its scale; the root and slope must then be constant.
\end{abstract}

\section{The quotient carrying the point data}

Use the curve $E/\Q(\sqrt2)$ and the $\Q$-simple abelian surface
$R=\operatorname{Res}_{\Q(\sqrt2)/\Q}E$ of \cite{Families}.
All point, quotient and trace functions and their guards are those
specified there. At a rational prime $p>7$ splitting in
$\Q(\sqrt2)$, the surface $R$ has good reduction: its two
local factors have integral equations with unit discriminant.

\begin{lemma}\label{lem:quotient}
Let $J/\Q$ be an abelian variety and let
$\Theta:J\to R^r$ be a morphism. Let $F\in\Q(J)$ be a
rational function depending only on $\Theta$. If $F=Y^2$ in
$\Q(J)$, the same square identity holds on an abelian variety
$J_1/\Q$ finite-isogenous to the translated image of $\Theta$.
The variety $J_1$ has good reduction at every good prime of $R$.

If a point-coordinate map $J_1\to R$ is nonconstant, its
reduction at every such prime is surjective. Consequently the
reciprocal abscissa of that coordinate has vertical valuation zero
in both slots at a split good prime.
\end{lemma}
\begin{proof}
Subtract $\Theta(0)$ to obtain a homomorphism $\theta$.
Put $J_1=J/\ker(\theta)^0$. The induced map from $J_1$ to
the abelian subvariety $S=\theta(J)\subset R^r$ is an isogeny.
The quotient $J\to J_1$ has geometrically connected fibers,
so $\Q(J_1)$ is algebraically closed in $\Q(J)$. Since $Y$
is algebraic over $\Q(J_1)$, it belongs to that field. This
also covers the zero square root.

For a prime $\ell$ distinct from the residue characteristic,
$V_\ell S$ is a subrepresentation of $V_\ell(R^r)$.
The Neron good-reduction criterion and isogeny invariance
\cite[Chapter IV, Theorem 3.5 and Corollary 3.6]{Milne} give
good reduction of $S$ and $J_1$ at every good prime of $R$.

A nonconstant point map has the form a translation followed by
a surjective homomorphism $\phi:J_1\to R$, by $\Q$-simplicity.
Poincare reducibility \cite[Chapter I, Proposition 10.1]{Milne}
supplies a homomorphism
$\psi:R\to J_1$ and a nonzero integer $m$ with
$\phi\psi=[m]$. These maps extend between the abelian schemes
at a good prime by the Neron mapping property
\cite[Chapter I, Remark 17.2]{Milne}.
Their reductions still satisfy this identity,
and $[m]$ is surjective even if the residue characteristic divides
$m$. Hence the reduction of $\phi$ is surjective; translation
preserves this conclusion.

The order at the generic point of the smooth special fiber is
a discrete valuation extending $v_p$ with $v(p)=1$.
Surjectivity of the reduced point map makes each reduced
abscissa a nonzero rational function at the generic point.
Thus $v(x)=v(1/x)=0$ in both slots.
\end{proof}

\section{A fixed-weight vertical obstruction}

Fix $X=x_0$, $A$ and a rational slope $z$ in the derivative
source. Thus $u,w$ are rational constants. Keep the actual functions
\[
 D_i=\frac{C_i^2}{1+uC_i^2},\quad
 W=\prod_i(1+D_i^2/u),\quad
 V=\Tr(W^{-1}),\quad w=\frac{u}{(1+u)^2},\quad
 \mathcal C=V^2+w(2-V).
\]
The point roles $H,U$ may vary independently.

\begin{theorem}\label{thm:derivative}
Suppose $p>7$ is split in $\Q(\sqrt2)$ and
$\nu=v_p(u)>0$ is odd. Every rational map from an abelian
variety over $\Q$ to this fixed-slope guarded square locus is
constant. The same conclusion holds for the three small-class
square tests of \cite{Families}.
\end{theorem}
\begin{proof}
Apply Lemma~\ref{lem:quotient} to the point map $(H,U)$ and
descend $Y$ to $J_1$. If the point map is constant, $z$ is fixed
and the ordinate has at most two values, so the whole map is constant.
Otherwise at least one of $H,U$ is a nonconstant map to $R$.
Its reciprocal is one of $C_1,C_2$ and has vertical valuation zero
in both slots by Lemma~\ref{lem:quotient}.

At the generic vertical valuation all nonzero tests have integral
valuations, possibly negative. The exact parity identities are
\[
 v(D_i)=\max(2v(C_i),-\nu),\qquad
 v(1+D_i^2/u)=\min(0,2v(D_i)-\nu)\le0.
\]
Zero tests contribute one. The selected unit reciprocal gives
factor valuation $-\nu$ in both slots. Thus
\[
 v(W_+^{-1}),v(W_-^{-1})\ge\nu,\qquad v(V)\ge\nu.
\]
Since $v(w)=\nu$ and $p\ne2$, the number $2-V$ is a unit.
The summands in $\mathcal C$ have valuations at least $2\nu$
and exactly $\nu$. Hence $v(\mathcal C)=\nu$ is odd,
contradicting the descended square identity. The three small
class multipliers have respective unit residues $2,2,1$, so
the same contradiction applies to them.
\end{proof}

\begin{theorem}\label{thm:general}
In the general two-trace source, fix $X=x_0$ and a rational
common slope $\lambda$. Suppose a split good prime $p>7$ satisfies
\[
                v_p(w_0)>0\text{ odd},\qquad
                v_p(\lambda+1)=0.
\]
Every rational map from an abelian variety over $\Q$ to its
guarded two-square locus is constant. The same conclusion holds
for the retained small-class alternatives.
\end{theorem}
\begin{proof}
The first comparison scale is one, so $w_0$ is constant.
Apply Lemma~\ref{lem:quotient} to all nine point roles and
descend both ordinates to $J_1$. If this point map is nonconstant,
at least one of
\[
                   H,\quad U_i,\quad U_i-Q_i\quad(1\le i\le4)
\]
is nonconstant: constancy of all $U_i$ and $U_i-Q_i$ would
force all $Q_i$ to be constant. Its reciprocal, or the conjugate
reciprocal, is one of the first ten field tests and is a unit in
both generic vertical slots.

Write $\nu=v_p(w_0)>0$ odd. Every field factor
$1+C_i^2/(\delta_0^2w_0^{2a_i-1})$ and every scalar factor
$1+G_{0,k}^2/w_0^{2\kappa_k-1}$ has nonpositive valuation by
parity. Zero tests contribute one. For the selected reciprocal
the field order is one, and $\delta_0=\lambda+1$ is a unit,
so its factor has valuation $-\nu$ in both slots. The preceding
trace argument gives $v(V_0)\ge\nu$ and
$v(\mathcal C_0)=\nu$ odd, a contradiction. Unit small-class
multipliers do not change it. A constant point map makes all
remaining finite ordinates constant as well.
\end{proof}

\begin{corollary}\label{cor:prepared}
For a fixed integer input and fixed root $A=[n]P$, a rational
map from an abelian variety over $\Q$ into the derivative
actual-trace locus whose image contains a prepared point is
constant. The same statement holds for the general two-trace
locus at a fixed input, with all nine point roles allowed to vary.
Finite point-chart covers and fixed menus of the stated small
classes do not change the assertion.
\end{corollary}
\begin{proof}
The nonconstant-slope obstructions of \cite{Families} first
make the slope constant. A prepared derivative point has selected
split good primes with $v_p(u)=1$. A prepared first-block general
contact has $v_p(w_0)=1$ and unit $\delta_0$. Thus the same
constant slope satisfies the respective fixed-weight hypothesis
throughout the family, and Theorem~\ref{thm:derivative} or
\ref{thm:general} applies. A curve or positive-dimensional
image cannot be contracted by a finite point-chart map.
\end{proof}

This excludes nonconstant abelian families through prepared points,
not the prepared points themselves. No Hasse principle or global
rational-witness assertion is used. A constant rational map exists
whenever an individual actual-trace witness exists. The conclusion
does not assert that all unprepared constant-slope loci are empty.

\section{Affine dependence of the derivative slope on its scale}

\begin{proposition}\label{prop:constant-root}
Fix $X=x_0$ and rational constants $m,b$. In any rational abelian
family in the guarded derivative square locus satisfying
$z=ma+b$, the root point $A$, its trace $T$, its scale $a$,
and its slope $z$ are constant. The same assertion holds for the
retained small-class alternatives. In particular it holds at a
fixed rational slope, by taking $m=0$.
This assertion does not require an odd valuation at a rational prime.
\end{proposition}
\begin{proof}
Before restricting to the line, the partial logarithmic derivative is
\[
 \frac{F_a}{F}=\sum_{c\in\mathcal S}
       \left(\frac1{a+c-z}+\frac1{a+c+z}\right).
\]
Suppose first that $m\ne\pm1$. The contact roots of the restricted
product form the two nine-element sets
\[
 \left\{\frac{b-c}{1-m}:c\in\mathcal S\right\},\qquad
 \left\{-\frac{b+c}{1+m}:c\in\mathcal S\right\}.
\]
Their union consists of $r\ge9$ distinct rational values
$a_1,\ldots,a_r$. At a root belonging only to the first or
second set, $u=-18F/F_a$ has derivative respectively
$-18(1-m)$ or $-18(1+m)$ with respect to $a$ along the line.
At a collision the two pole residues in $F_a/F$ add to
\[
             \frac1{1-m}+\frac1{1+m}=\frac2{1-m^2}\ne0.
\]
Consequently the restricted derivative of $u$ is
$-9(1-m^2)$ there. Thus $u$ has a simple zero at every $a_j$,
including every collision.

If $m=1$, one family of factors is constant. A zero constant
factor would make $F$ identically zero and violates the retained
guards. Otherwise the second family gives nine distinct rational
roots, with derivative $-36$ for $u$ at each. The case $m=-1$
is identical with the two families interchanged. In all guarded
cases, $u$ therefore has simple zeros at at least nine distinct
rational values of $a$.

The geometric divisor argument of \cite{Families} uses only the
two reference tests $C_1,C_3$: their simultaneous zero loci in a
slot pull back from a codimension-two point of $R$, unless $H$
is constant, when both are nonzero constants. Accordingly any
divisor with positive odd order of $u$ would give odd order of
the radicand. A square identity therefore forces every zero
divisor of $a-a_j$ to be even.

If $a$ were nonconstant, the product
$q=\prod_{j=1}^r(a-a_j)$ would have even divisor. Its zero
divisors are even by the preceding paragraph; its pole divisor
is even because $a=M^4$ has all pole multiplicities divisible by
four, even when $r$ is odd. A square root of $q$ on a finite etale
abelian cover would give a nonconstant map to
\[
                         s^2=\prod_{j=1}^r(t-a_j),
\]
whose smooth normalization has genus
$\lfloor(r-1)/2\rfloor\ge4$. The Jacobian argument in
\cite{Families} forbids this, so $a$ is constant.

Since $a=((x_0^2+B)^4+T^4)^4$, the function $T$ is algebraic
over $\Q$, hence constant on the geometrically integral abelian
variety. A nonconstant $A$-map would surject onto the $\Q$-simple
surface $R$, but $\Tr x$ is nonconstant on that surface: after
base change it is the sum of the two independent abscissa functions.
Thus $A$ is constant as well, and the affine relation makes $z$
constant. The small-class multipliers are units at the geometric
odd-valuation contradiction, so they do not change the proof.
\end{proof}

This proposition treats fixed slopes and rational affine dependence
on the scale. Unrestricted simultaneous variation of the derivative
root and slope is outside its conclusion. The actual fixed-root
realization criterion is already covered by
Corollary~\ref{cor:prepared}. No defining-variable reduction follows
from these obstructions alone.

\begin{thebibliography}{9}
\bibitem{Families}\emph{Curves and abelian families in an actual-trace square locus},
\nolinkurl{negative_results/papers/pole_parameters/curves_and_abelian_families_20261117_v69.tex}.
\bibitem{Milne} J. S. Milne, \emph{Abelian Varieties}, version 2.0,
2008. Chapter I, Proposition 10.1 and Remark 17.2, printed
pages 42 and 70, give reducibility and the Neron mapping property.
Chapter IV, Theorem 3.5 and Corollary 3.6, printed
pages 141--142, give the good-reduction criterion and isogeny invariance.
\url{https://www.jmilne.org/math/CourseNotes/AV.pdf}.
\end{thebibliography}
\end{document}
